\documentclass[12pt]{article}
% ----------------------------------------------------------------------
% Define external packages, language, margins, fonts, new commands 
% and colors
% ----------------------------------------------------------------------
\usepackage[utf8]{inputenc} % Codification
\usepackage[english]{babel} % Writing idiom

\usepackage[export]{adjustbox} % Align images
\usepackage{amsmath} % Extra commands for math mode
\usepackage{amssymb} % Mathematical symbols
\usepackage{anysize} % Personalize margins
    \marginsize{2cm}{2cm}{2cm}{2cm} % {left}{right}{above}{below}
\usepackage{appendix} % Appendices
\usepackage{cancel} % Expression cancellation
\usepackage{caption} % Captions
    \captionsetup{labelfont={bf}}
\usepackage{cite} % Citations, like [1 - 3]
\usepackage{color} % Text coloring
\usepackage{fancyhdr} % Head note and footnote
    \pagestyle{fancy}
    \fancyhf{}
    \fancyhead[L]{
        \includegraphics[scale = 0.4]{NovaFct.png}
    } % Left of Head note
    \fancyhead[R]{\footnotesize Eletrónica para Micro-Sistemas} % Right of Head note
    \fancyfoot[L]{\footnotesize MEEC/MIEEC} % Left of Footnote
    \fancyfoot[R]{\thepage  } % Center of Footnote
    %\fancyfoot[R]{\footnotesize Degree} % Right of Footnote
    \renewcommand{\footrulewidth}{0.4pt} % Footnote rule
\usepackage{float} % Utilization of [H] in figures
\usepackage{graphicx} % Figures in LaTeX
\usepackage[colorlinks = true, plainpages = true, linkcolor = novablue, urlcolor = novablue, citecolor = novablue, anchorcolor = novablue]{hyperref}
\usepackage{indentfirst} % First paragraph
\usepackage[super]{nth} % Superscripts
\usepackage{siunitx} % SI units
\usepackage{subcaption} % Subfigures
\usepackage{titlesec} % Font
    \titleformat{\section}{\Large\bfseries}{\thesection}{1em}{}
    \titleformat{\subsection}{\large\bfseries}{\thesubsection}{1em}{}
    \titleformat{\subsubsection}{\normalsize\bfseries}{\thesubsubsection}{1em}{}
    %\fancyfoot[C]{\thepage}

%code
\usepackage[breakable]{tcolorbox}

% Code highlighting
\usepackage{listings}
\lstset{ 
    backgroundcolor=\color{cellbackground}, % background color for the code block
    basicstyle=\ttfamily,                   % font style
    breakatwhitespace=false,                % automatic breaks only at whitespace
    breaklines=true,                        % automatic line breaking
    captionpos=b,                           % caption position
    commentstyle=\color{gray},              % comment style
    escapeinside={\%*}{*)},                 % if you want to add LaTeX within your code
    keywordstyle=\color{blue},              % keyword style
    stringstyle=\color{dkgreen},            % string style
    frame=single,                           % adds a frame around the code
    rulecolor=\color{cellborder},           % border color
}

% Exact colors from NB
\definecolor{incolor}{HTML}{303F9F}
\definecolor{outcolor}{HTML}{D84315}
\definecolor{cellborder}{HTML}{CFCFCF}
\definecolor{cellbackground}{HTML}{F7F7F7}

% Random text (not needed)
\usepackage{lipsum}
\usepackage{duckuments}

% New and re-newcommands
\newcommand{\sen}{\operatorname{\sen}} % Sine function definition
\newcommand{\HRule}{\rule{\linewidth}{0.5mm}} % Specific rule definition
\renewcommand{\appendixpagename}{\LARGE Appendices}

% Colors
\definecolor{novablue}{RGB}{0, 101, 189}
\definecolor{dkgreen}{rgb}{0,0.6,0} 

%%%%%%%%%%%%%%%%%%%%%%%%%%%%%%%%%%%%%%%%%%%%%%%%%%%%%%%%%%%%%%%%%%%%%%%%
%                                 Document                             %
%%%%%%%%%%%%%%%%%%%%%%%%%%%%%%%%%%%%%%%%%%%%%%%%%%%%%%%%%%%%%%%%%%%%%%%%
\begin{document}

% ----------------------------------------------------------------------
% Cover
% ----------------------------------------------------------------------
\begin{center}
    \begin{figure}
        \vspace{-1.0cm}
        \includegraphics[scale = 1, left]{NovaFct.png} % Nova logo
    \end{figure}

    \mbox{}\\[2.0cm]
    \textsc{\Huge MEEC/MIEEC}\\[2.5cm]
    \textsc{\LARGE Eletrónica para Micro-Sistemas}\\[2.0cm]
    \HRule\\[0.4cm]
    {\large \bf { Laboratório 1 Sensor de Temperatura } [\texttt{PT}]}\\[0.2cm]
    \HRule\\[1.5cm]
\end{center}

\begin{flushleft}
    \textbf{Authors:}
\end{flushleft}

\begin{center}
    \begin{minipage}{0.5\textwidth}
        \begin{flushleft}
            Martim Duarte Agostinho (70392)\\
            Lorem Ipsum (\texttt{ISTID} $\in \mathbb{Z}^+$)\\
            Lorem Ipsum (\texttt{ISTID} $\in \mathbb{Z}^+$)
        \end{flushleft}
    \end{minipage}%
    \begin{minipage}{0.5\textwidth}
        \begin{flushright}
            \href{mailto:md.agostinho@campus.fct.unl.pt}{\texttt{md.agostinho@campus.fct.unl.pt}}\\
            \href{mailto:lorem.ipsum@campus.fct.unl.pt}{\texttt{lorem.ipsum@campus.fct.unl.pt}}\\
            \href{mailto:lorem.ipsum@campus.fct.unl.pt}{\texttt{lorem.ipsum@campus.fct.unl.pt}}
        \end{flushright}
    \end{minipage}
\end{center}
 
% Caixa a dizer o grupo
%\begin{flushleft}
%    \large $\boxed{\text{\bf Group} \ \clubsuit}$\\[4.0cm]
%\end{flushleft}

\vspace{6cm}

\begin{center}
    \large \bf 2023/2024 -- \nth{2} Semester, P4
\end{center}

\thispagestyle{empty}

\setcounter{page}{0}

\newpage

\newpage

\tableofcontents % Generates the table of contents

\newpage

\listoffigures

\newpage

\section{Dimensionamento}
\subsection{NTC}

Para puder dimensionar o AFE do NTC primeiro é necessario definir o intrevalo de temperaturas em que o circuito irá operar.
Foi então decidido que seria adequado um temperatura $T \in [ 10^{\circ}; 40^{\circ} ]$.
E pelo datasheet do NTC foi obtido o intrevalo da sua resistencia $R_{NTC} \in [~5,282k~;~19,98k~]$   

Para usar equacao Steinhart-Hart $\frac{1}{T} = A + B\cdot \ln(R) + C\cdot [\ln(R)]^3$, precisamos de usar 3 pontos para encontrar as constantes $A$, $B$ e $C$.
~
$R(T) = R_{NTC} $ onde $T$ é a temperatura em kelvin e $R_{NTC}$ é o valor da resistencia do thermistor \textit{NTC}

\begin{equation}
  \left\{\begin{array}{@{}l@{}}
    R( 283,15 ) = 1,998\cdot 10^4 ~\Omega \\
    R( 298,15 ) = 10^4 ~\Omega\\
    R( 313,15 ) = 0,5282 \cdot 10^4 ~\Omega\\
  
    \end{array}\right.\
\end{equation}

\begin{equation}
  \left\{\begin{array}{@{}l@{}}
     A = 1,2 \cdot 10^{-3}\\
     B = 2,1 \cdot 10^{-4}\\
     C = 1,3 \cdot 10^{-7}\\
  
    \end{array}\right.\
\end{equation}



\subsection{LM35}
\subsection{DS18B20 ??(Este tem Dimensionamento?)}
\subsection{Relé de saída }


\section{ Simulação }
\section{ Implementação }

    \begin{figure}[H]
        \centering
        \includegraphics*[scale = 0.5]{NovaFct.png}
        \caption{Logo da Nova FCT}
        \label{wrap-fig:1}
    \end{figure}

\end{document}